\documentclass[11pt]{article}
\usepackage[utf8]{inputenc}
\usepackage[spanish]{babel}
\usepackage{geometry}
\usepackage{booktabs}
\usepackage{longtable}
\usepackage{array}
\usepackage{xcolor}
\usepackage{hyperref}
\usepackage{fancyvrb}

\geometry{a4paper, landscape, margin=1cm}

\title{Documentación de Endpoints de la API - K-Pop Groups Manager}
\author{IS-K Pop Backend}
\date{15 de enero de 2026}

\begin{document}

\maketitle

\section*{Información General}

\textbf{Base URL:} \texttt{http://localhost:\{PORT\}/api}

\textbf{Nota:} Todos los DTOs mostrados son los datos que el frontend debe enviar al backend en el body de las peticiones.

\section{Autenticación (\texttt{/auth})}

\begin{longtable}{lllll}
\toprule
\textbf{Método} & \textbf{URL} & \textbf{Autenticación} & \textbf{Rol Requerido} & \textbf{Descripción} \\
\midrule
\endhead
POST & \texttt{/api/auth/login} & No & - & Iniciar sesión \\
\bottomrule
\end{longtable}

\textbf{DTO Login:}
\begin{Verbatim}
{
  "email": "string",
  "password": "string"
}
\end{Verbatim}

\section{Usuarios (\texttt{/user})}

\begin{longtable}{lllll}
\toprule
\textbf{Método} & \textbf{URL} & \textbf{Autenticación} & \textbf{Rol Requerido} & \textbf{Descripción} \\
\midrule
\endhead
POST & \texttt{/api/user} & No & - & Crear usuario \\
GET & \texttt{/api/user} & No & - & Listar todos los usuarios \\
GET & \texttt{/api/user/:id} & No & - & Obtener usuario por ID \\
PUT & \texttt{/api/user/:id} & No & - & Actualizar usuario \\
DELETE & \texttt{/api/user/:id} & No & - & Eliminar usuario \\
\bottomrule
\end{longtable}

\textbf{DTO Crear Usuario (POST /api/user):}
\begin{Verbatim}
{
  "email": "string",
  "name": "string",
  "password": "string",
  "role": "ADMIN | STAFF | ARTIST | APPRENTICE | MANAGER | DIRECTOR",
  "agencyId": "number (opcional)",
  "IdAp": "number (opcional)",
  "IdGr": "number (opcional)"
}
\end{Verbatim}

\textbf{DTO Actualizar Usuario (PUT /api/user/:id):}
\begin{Verbatim}
{
  "email": "string (opcional)",
  "name": "string (opcional)",
  "role": "string (opcional)"
}
\end{Verbatim}

\section{Agencias (\texttt{/agency})}

\begin{longtable}{lllll}
\toprule
\textbf{Método} & \textbf{URL} & \textbf{Autenticación} & \textbf{Rol Requerido} & \textbf{Descripción} \\
\midrule
\endhead
POST & \texttt{/api/agency} & Sí & Staff & Crear agencia \\
GET & \texttt{/api/agency} & No & - & Listar todas las agencias \\
GET & \texttt{/api/agency/:id} & No & - & Obtener agencia por ID \\
PUT & \texttt{/api/agency/:id} & Sí & Staff & Actualizar agencia \\
DELETE & \texttt{/api/agency/:id} & Sí & Staff & Eliminar agencia \\
GET & \texttt{/api/agency/search/agency\_name} & No & - & Buscar agencias por nombre \\
GET & \texttt{/api/agency/search/agency\_address} & No & - & Buscar agencias por dirección \\
GET & \texttt{/api/agency/search/agency\_foundation} & No & - & Buscar por fecha fundación \\
POST & \texttt{/api/agency/:idAgency/artists} & No & - & Obtener artistas con contratos activos \\
POST & \texttt{/api/agency/:idAgency/groups} & No & - & Obtener grupos con contratos activos \\
\bottomrule
\end{longtable}

\textbf{DTO Crear Agencia (POST /api/agency):}
\begin{Verbatim}
{
  "name": "string",
  "address": "string",
  "foundation": "date (YYYY-MM-DD)"
}
\end{Verbatim}

\section{Aprendices (\texttt{/apprentice})}

\begin{longtable}{lllll}
\toprule
\textbf{Método} & \textbf{URL} & \textbf{Autenticación} & \textbf{Rol Requerido} & \textbf{Descripción} \\
\midrule
\endhead
POST & \texttt{/api/apprentice/:id} & Sí & Staff & Crear aprendiz (id = agencyId) \\
GET & \texttt{/api/apprentice} & Sí & - & Listar todos los aprendices \\
GET & \texttt{/api/apprentice/:id} & Sí & - & Obtener aprendiz por ID \\
GET & \texttt{/api/apprentice/name/:name} & Sí & - & Obtener aprendiz por nombre \\
GET & \texttt{/api/apprentice/agency/:id} & Sí & - & Listar aprendices por agencia \\
GET & \texttt{/api/apprentice/:id/evaluations} & Sí & - & Obtener evaluaciones del aprendiz \\
PUT & \texttt{/api/apprentice/:id} & Sí & Staff & Actualizar aprendiz \\
PUT & \texttt{/api/apprentice/attract/:apprenticeId/:agencyId} & Sí & - & Atraer aprendiz a agencia \\
POST & \texttt{/api/apprentice/evaluation} & Sí & - & Añadir evaluación a aprendiz \\
GET & \texttt{/api/apprentice/scout} & Sí & - & Obtener aprendices disponibles para scout \\
DELETE & \texttt{/api/apprentice/:id} & Sí & Staff & Eliminar aprendiz \\
\bottomrule
\end{longtable}

\textbf{DTO Crear Aprendiz (POST /api/apprentice/:id):}
\begin{Verbatim}
{
  "name": "string",
  "dateOfBirth": "date (YYYY-MM-DD)",
  "agencyId": "number (obtenido de la ruta)"
}
Nota: La edad debe estar entre 15 y 50 años
\end{Verbatim}

\textbf{DTO Evaluación (POST /api/apprentice/evaluation):}
\begin{Verbatim}
{
  "apprenticeId": "number",
  "score": "number",
  "evaluationDate": "date",
  "comments": "string (opcional)"
}
\end{Verbatim}

\section{Artistas (\texttt{/artist})}

\begin{longtable}{lllll}
\toprule
\textbf{Método} & \textbf{URL} & \textbf{Autenticación} & \textbf{Rol Requerido} & \textbf{Descripción} \\
\midrule
\endhead
POST & \texttt{/api/artist} & Sí & Staff & Crear artista \\
GET & \texttt{/api/artist} & Sí & - & Listar todos los artistas \\
GET & \texttt{/api/artist/:id} & Sí & - & Obtener artistas por agencia (id = agencyId) \\
GET & \texttt{/api/artist/:apprenticeId\&:groupId} & Sí & - & Obtener artista específico por IDs compuestos \\
GET & \texttt{/api/artist/query/:id} & Sí & - & Obtener artistas en debut \\
POST & \texttt{/api/artist/income/} & Sí & - & Obtener ingresos totales del artista \\
POST & \texttt{/api/artist/bestAlbums/} & Sí & - & Obtener mejores álbumes del artista \\
POST & \texttt{/api/artist/succes/} & Sí & - & Obtener ingresos y éxito del artista \\
GET & \texttt{/api/artist/agencyChanges/} & Sí & - & Obtener cambios de agencia \\
GET & \texttt{/api/artist/soloArtists/} & Sí & - & Obtener historial profesional de artistas solo \\
PUT & \texttt{/api/artist/:apprenticeId\&:groupId} & Sí & Staff & Actualizar artista \\
DELETE & \texttt{/api/artist/:apprenticeId\&:groupId} & Sí & Staff & Eliminar artista \\
\bottomrule
\end{longtable}

\textbf{DTO Crear Artista (POST /api/artist):}
\begin{Verbatim}
{
  "ApprenticeId": "number",
  "GroupId": "number",
  "ArtistName": "string",
  "DebutDate": "date (YYYY-MM-DD)",
  "Status": "ACTIVO | INACTIVO | RETIRADO"
}
\end{Verbatim}

\section{Álbumes (\texttt{/album})}

\begin{longtable}{lllll}
\toprule
\textbf{Método} & \textbf{URL} & \textbf{Autenticación} & \textbf{Rol Requerido} & \textbf{Descripción} \\
\midrule
\endhead
POST & \texttt{/api/album} & No & - & Crear álbum \\
GET & \texttt{/api/album} & No & - & Listar todos los álbumes \\
GET & \texttt{/api/album/:id} & No & - & Obtener álbum por ID \\
PUT & \texttt{/api/album/:id} & No & - & Actualizar álbum \\
DELETE & \texttt{/api/album/:id} & No & - & Eliminar álbum \\
\bottomrule
\end{longtable}

\textbf{DTO Crear Álbum (POST /api/album):}
\begin{Verbatim}
{
  "apprenticeId": "number",
  "groupId": "number",
  "title": "string",
  "releaseDate": "date (YYYY-MM-DD)",
  "producer": "string"
}
\end{Verbatim}

\section{Canciones (\texttt{/song})}

\begin{longtable}{lllll}
\toprule
\textbf{Método} & \textbf{URL} & \textbf{Autenticación} & \textbf{Rol Requerido} & \textbf{Descripción} \\
\midrule
\endhead
POST & \texttt{/api/song} & No & - & Crear canción \\
GET & \texttt{/api/song} & No & - & Listar todas las canciones \\
GET & \texttt{/api/song/:id} & No & - & Obtener canción por ID \\
PUT & \texttt{/api/song/:id} & No & - & Actualizar canción \\
DELETE & \texttt{/api/song/:id} & No & - & Eliminar canción \\
\bottomrule
\end{longtable}

\textbf{DTO Crear Canción (POST /api/song):}
\begin{Verbatim}
{
  "title": "string",
  "gender": "string (género musical)",
  "producer": "string",
  "releaseDate": "date (YYYY-MM-DD)",
  "albumIds": ["number, number, ..."] (opcional)
}
\end{Verbatim}

\section{Aplicaciones (\texttt{/application})}

\begin{longtable}{lllll}
\toprule
\textbf{Método} & \textbf{URL} & \textbf{Autenticación} & \textbf{Rol Requerido} & \textbf{Descripción} \\
\midrule
\endhead
POST & \texttt{/api/application} & No & - & Crear aplicación \\
POST & \texttt{/api/application/create} & No & - & Crear aplicación con rol \\
GET & \texttt{/api/application} & No & - & Listar todas las aplicaciones \\
GET & \texttt{/api/application/:id} & No & - & Obtener aplicación por ID \\
GET & \texttt{/api/application/apprentices/without-application} & No & - & Obtener aprendices sin aplicación \\
GET & \texttt{/api/application/soloistArtist/} & No & - & Obtener artistas solista \\
GET & \texttt{/api/application/createGroup/:id} & No & - & Crear grupo a partir de aplicación \\
PUT & \texttt{/api/application/:id} & No & - & Actualizar aplicación \\
PATCH & \texttt{/api/application/:id/apprentice-decision} & No & - & Decisión del aprendiz en aplicación \\
PATCH & \texttt{/api/application/:id/artist-decision} & No & - & Decisión del artista en aplicación \\
DELETE & \texttt{/api/application/:id} & No & - & Eliminar aplicación \\
\bottomrule
\end{longtable}

\textbf{DTO Crear Aplicación (POST /api/application):}
\begin{Verbatim}
{
  "groupName": "string",
  "idAgency": "number",
  "idConcept": "number",
  "apprentices": [["number", "string"]],
  "artists": [["number", "number", "string"]],
  "status": "string (PENDIENTE por defecto)"
}
\end{Verbatim}

\section{Conceptos (\texttt{/concept})}

\begin{longtable}{lllll}
\toprule
\textbf{Método} & \textbf{URL} & \textbf{Autenticación} & \textbf{Rol Requerido} & \textbf{Descripción} \\
\midrule
\endhead
POST & \texttt{/api/concept} & No & - & Crear concepto \\
GET & \texttt{/api/concept} & No & - & Listar todos los conceptos \\
GET & \texttt{/api/concept/:id} & No & - & Obtener concepto por ID \\
PUT & \texttt{/api/concept/:id} & No & - & Actualizar concepto \\
DELETE & \texttt{/api/concept/:id} & No & - & Eliminar concepto \\
\bottomrule
\end{longtable}

\textbf{DTO Crear Concepto (POST /api/concept):}
\begin{Verbatim}
{
  "name": "string",
  "description": "string"
}
\end{Verbatim}

\section{Grupos (\texttt{/group})}

\begin{longtable}{lllll}
\toprule
\textbf{Método} & \textbf{URL} & \textbf{Autenticación} & \textbf{Rol Requerido} & \textbf{Descripción} \\
\midrule
\endhead
POST & \texttt{/api/group} & No & - & Crear grupo \\
GET & \texttt{/api/group} & No & - & Listar todos los grupos \\
GET & \texttt{/api/group/:id} & No & - & Obtener grupo por ID \\
GET & \texttt{/api/group/search/name} & No & - & Buscar grupo por nombre \\
GET & \texttt{/api/group/search/debut} & No & - & Buscar grupos por fecha debut \\
GET & \texttt{/api/group/search/status} & No & - & Buscar grupos por estado \\
GET & \texttt{/api/group/search/membercount} & No & - & Buscar grupos por cantidad miembros \\
GET & \texttt{/api/group/search/member} & No & - & Buscar grupos por miembro \\
GET & \texttt{/api/group/search/agency} & No & - & Buscar grupos por agencia \\
GET & \texttt{/api/group/search/concept} & No & - & Buscar grupos por concepto \\
GET & \texttt{/api/group/search/visualconcept} & No & - & Buscar grupos por concepto visual \\
PUT & \texttt{/api/group/update/:id} & No & - & Actualizar grupo \\
DELETE & \texttt{/api/group/delete/:id} & No & - & Eliminar grupo \\
POST & \texttt{/api/group/:id/members/add} & No & - & Agregar miembros al grupo \\
POST & \texttt{/api/group/:id/members/remove} & No & - & Remover miembros del grupo \\
POST & \texttt{/api/group/:id/albums/add} & No & - & Agregar álbumes al grupo \\
POST & \texttt{/api/group/:id/activities/add} & No & - & Agregar actividades al grupo \\
\bottomrule
\end{longtable}

\textbf{DTO Crear Grupo (POST /api/group):}
\begin{Verbatim}
{
  "name": "string",
  "debut": "date (YYYY-MM-DD)",
  "status": "ACTIVO | INACTIVO | DISUELTO",
  "memberCount": "number",
  "IdAgency": "number",
  "IdConcept": "number",
  "IdVisualConcept": "number",
  "members": ["number, number, ..."],
  "roles": ["string, string, ..."],
  "albums": ["number, number, ..."] (opcional),
  "activities": ["number, number, ..."] (opcional)
}
Nota: members y roles deben tener la misma longitud
\end{Verbatim}

\section{Actividades (\texttt{/activity})}

\begin{longtable}{lllll}
\toprule
\textbf{Método} & \textbf{URL} & \textbf{Autenticación} & \textbf{Rol Requerido} & \textbf{Descripción} \\
\midrule
\endhead
POST & \texttt{/api/activity} & Sí & - & Crear actividad \\
GET & \texttt{/api/activity} & Sí & - & Listar todas las actividades \\
GET & \texttt{/api/activity/:id} & Sí & - & Obtener actividad por ID \\
GET & \texttt{/api/activity/:apprenticeId\&:groupId} & Sí & - & Obtener actividades por artista \\
GET & \texttt{/api/activity/byGroup/:groupId} & Sí & - & Obtener actividades por grupo \\
PUT & \texttt{/api/activity/:id} & Sí & Staff & Actualizar actividad \\
PUT & \texttt{/api/activity/:id/decision} & Sí & - & Tomar decisión sobre actividad \\
DELETE & \texttt{/api/activity/:id} & Sí & Staff & Eliminar actividad \\
POST & \texttt{/api/activity/addArtist} & Sí & Staff & Agregar artista a actividad \\
\bottomrule
\end{longtable}

\textbf{DTO Crear Actividad (POST /api/activity):}
\begin{Verbatim}
{
  "responsible": "string",
  "activityType": "CONCIERTO | PROMOCION | EVENTO | OTRO",
  "date": "date (YYYY-MM-DD)",
  "place": "string",
  "eventType": "string",
  "artists": [["number", "number"], ...] (opcional),
  "groups": ["number", ...] (opcional)
}
\end{Verbatim}

\section{Contratos (\texttt{/contract})}

\begin{longtable}{lllll}
\toprule
\textbf{Método} & \textbf{URL} & \textbf{Autenticación} & \textbf{Rol Requerido} & \textbf{Descripción} \\
\midrule
\endhead
POST & \texttt{/api/contract} & No & - & Crear contrato \\
GET & \texttt{/api/contract} & No & - & Listar todos los contratos \\
GET & \texttt{/api/contract/find} & No & - & Obtener contrato por ID (query param) \\
GET & \texttt{/api/contract/offer/artist} & No & - & Obtener ofertas de contrato artista \\
GET & \texttt{/api/contract/offer/group} & No & - & Obtener ofertas de contrato grupo \\
PUT & \texttt{/api/contract} & No & - & Actualizar contrato \\
DELETE & \texttt{/api/contract} & No & - & Eliminar contrato \\
\bottomrule
\end{longtable}

\textbf{DTO Crear Contrato (POST /api/contract):}
\begin{Verbatim}
{
  "type": "ARTISTA | APRENDIZ | GRUPO",
  "agencyId": "number",
  "startDate": "date (YYYY-MM-DD)",
  "status": "string (Pendiente por defecto)",
  "initialConditions": "string",
  "incomeDistribution": "string",
  "apprenticeId": "number (opcional, para contratos artista)",
  "groupId": "number (opcional, para contratos grupo)",
  "completionDate": "date (opcional)"
}
\end{Verbatim}

\section{Premios (\texttt{/award})}

\begin{longtable}{lllll}
\toprule
\textbf{Método} & \textbf{URL} & \textbf{Autenticación} & \textbf{Rol Requerido} & \textbf{Descripción} \\
\midrule
\endhead
POST & \texttt{/api/award} & No & - & Crear premio \\
GET & \texttt{/api/award} & No & - & Listar todos los premios \\
GET & \texttt{/api/award/:id} & No & - & Obtener premio por ID \\
PUT & \texttt{/api/award/:id} & No & - & Actualizar premio \\
DELETE & \texttt{/api/award/:id} & No & - & Eliminar premio \\
\bottomrule
\end{longtable}

\textbf{DTO Crear Premio (POST /api/award):}
\begin{Verbatim}
{
  "awardTitle": "string",
  "academyName": "string"
}
\end{Verbatim}

\section{Ingresos (\texttt{/income})}

\begin{longtable}{lllll}
\toprule
\textbf{Método} & \textbf{URL} & \textbf{Autenticación} & \textbf{Rol Requerido} & \textbf{Descripción} \\
\midrule
\endhead
POST & \texttt{/api/income} & No & - & Crear ingreso \\
GET & \texttt{/api/income} & No & - & Listar todos los ingresos \\
GET & \texttt{/api/income/:id} & No & - & Obtener ingreso por ID \\
PUT & \texttt{/api/income/:id} & No & - & Actualizar ingreso \\
DELETE & \texttt{/api/income/:id} & No & - & Eliminar ingreso \\
\bottomrule
\end{longtable}

\textbf{DTO Crear Ingreso (POST /api/income):}
\begin{Verbatim}
{
  "description": "string",
  "amount": "number",
  "idActivity": "number"
}
\end{Verbatim}

\section{Listas de Popularidad (\texttt{/popularity})}

\begin{longtable}{lllll}
\toprule
\textbf{Método} & \textbf{URL} & \textbf{Autenticación} & \textbf{Rol Requerido} & \textbf{Descripción} \\
\midrule
\endhead
POST & \texttt{/api/popularity} & No & - & Crear lista de popularidad \\
GET & \texttt{/api/popularity} & No & - & Listar todas las listas \\
GET & \texttt{/api/popularity/:id} & No & - & Obtener lista por ID \\
PUT & \texttt{/api/popularity/:id} & No & - & Actualizar lista \\
DELETE & \texttt{/api/popularity/:id} & No & - & Eliminar lista \\
POST & \texttt{/api/popularity/addSong} & No & - & Agregar canción a lista \\
POST & \texttt{/api/popularity/delSong} & No & - & Eliminar canción de lista \\
POST & \texttt{/api/popularity/ranking} & No & - & Actualizar ranking de lista \\
\bottomrule
\end{longtable}

\textbf{DTO Crear Lista de Popularidad (POST /api/popularity):}
\begin{Verbatim}
{
  "name": "string",
  "listType": "Nacional | Internacional",
  "requirement": "number (opcional)"
}
\end{Verbatim}

\section{Conceptos Visuales (\texttt{/visual-concept})}

\begin{longtable}{lllll}
\toprule
\textbf{Método} & \textbf{URL} & \textbf{Autenticación} & \textbf{Rol Requerido} & \textbf{Descripción} \\
\midrule
\endhead
GET & \texttt{/api/visual-concept} & No & - & Listar todos los conceptos visuales \\
GET & \texttt{/api/visual-concept/:id} & No & - & Obtener concepto visual por ID \\
DELETE & \texttt{/api/visual-concept/:id} & No & - & Eliminar concepto visual \\
\bottomrule
\end{longtable}

\textbf{DTO Crear Concepto Visual (POST /api/visual-concept):}
\begin{Verbatim}
{
  "picture": "string (ruta o URL de imagen)"
}
\end{Verbatim}

\section{Notas Importantes}

\subsection*{Autenticación}
Los endpoints que requieren autenticación necesitan un token JWT en el header:

\texttt{Authorization: Bearer \{token\}}

\subsection*{Roles Disponibles}
\begin{itemize}
    \item \textbf{Admin}: Acceso completo al sistema - puede gestionar todos los usuarios
    \item \textbf{Staff}: Director/Manager - puede gestionar agencias, aprendices, artistas y actividades
    \item \textbf{Agency Access}: Rol especial para acceder a recursos de agencias específicas (requireAgencyAccess middleware)
    \item \textbf{Artist}: Artista del sistema
    \item \textbf{Apprentice}: Aprendiz del sistema
    \item \textbf{Manager}: Manager de una agencia
    \item \textbf{Director}: Director de una agencia
\end{itemize}

\subsection*{Middlewares de Autenticación}
\begin{itemize}
    \item \textbf{AuthMiddleware.authenticate()}: Verifica token JWT válido
    \item \textbf{RoleMiddleware.onlyAdmin()}: Solo permite acceso a usuarios Admin
    \item \textbf{RoleMiddleware.onlyStaff()}: Permite acceso a Director y Manager
    \item \textbf{RoleMiddleware.requireAgencyAccess}: Verifica acceso a recursos de la agencia
    \item \textbf{RoleMiddleware.updatePropertiesId()}: Permite actualizar si eres dueño del recurso o Admin
\end{itemize}

\section{Notas Importantes}

\subsection*{Autenticación}
Los endpoints que requieren autenticación necesitan un token JWT en el header:

\texttt{Authorization: Bearer \{token\}}

\subsection*{Roles Disponibles}
\begin{itemize}
    \item \textbf{ADMIN}: Acceso completo al sistema - puede gestionar todos los usuarios
    \item \textbf{STAFF}: Director/Manager - puede gestionar agencias, aprendices, artistas y actividades
    \item \textbf{ARTIST}: Artista del sistema
    \item \textbf{APPRENTICE}: Aprendiz del sistema
    \item \textbf{MANAGER}: Manager de una agencia
    \item \textbf{DIRECTOR}: Director de una agencia
\end{itemize}

\subsection*{Middlewares de Autenticación}
\begin{itemize}
    \item \textbf{AuthMiddleware.authenticate()}: Verifica token JWT válido
    \item \textbf{RoleMiddleware.onlyStaff()}: Permite acceso a Director y Manager
    \item \textbf{RoleMiddleware.updatePropertiesId()}: Permite actualizar si eres dueño del recurso o Admin
\end{itemize}

\subsection*{Estructura de DTOs}

\textbf{Importante:} Todos los DTOs deben estar en el formato JSON y enviarse en el body de la petición, excepto para operaciones GET que usan query parameters.

Los siguientes tipos de datos son comunes:
\begin{itemize}
    \item \textbf{date}: Formato ISO 8601 (YYYY-MM-DD)
    \item \textbf{number}: Valores numéricos (enteros o decimales)
    \item \textbf{string}: Texto
    \item \textbf{boolean}: true/false
    \item \textbf{arrays}: Listas entre corchetes []
\end{itemize}

\subsection*{Notas de Implementación}
\begin{itemize}
    \item Los endpoints de usuario NO requieren autenticación actualmente (comentada en UserRoutes)
    \item Los endpoints de grupos y conceptos no requieren autenticación
    \item Los endpoints de artista usan IDs compuestos: \texttt{:apprenticeId\&:groupId} (clave primaria compuesta)
    \item Los endpoints de contratos no usan parámetros en la ruta para PUT/DELETE, se envía el id en el body
    \item El endpoint \texttt{GET /api/contract/find} espera el id como query parameter
    \item El endpoint \texttt{POST /api/apprentice/:id} recibe el agencyId en la ruta
    \item La edad de un aprendiz debe estar entre 15 y 50 años
    \item Para crear un grupo, la cantidad de miembros debe corresponder con la longitud del array de miembros
    \item El status de aprendiz se define automáticamente al crear
    \item Para actividades, los artistas se envían como pares [apprenticeId, groupId]
\end{itemize}

\subsection*{Configuración}
El puerto debe ser configurado en el archivo \texttt{secrets.ts}.

\subsection*{Últimas Actualizaciones (15 de enero de 2026)}
\begin{itemize}
    \item Documentación completa de DTOs agregada para todos los endpoints
    \item Endpoints de Album, Song, Activity, etc. actualizados con información correcta
    \item Se agregó sección de Ingresos (\texttt{/income})
    \item Se aclararon endpoints de búsqueda en grupos
    \item Se incluye información sobre estructuras de datos esperadas en cada endpoint
\end{itemize}

\end{document}
